% Transcription — fonds Alexandre Grothendieck, shelfmark 60, batch 2 (pages 21-40).
% Established from the facsimile archives/batches/60/batch-02.pdf (PDF page k =
% archivists' page 20+k, no cover sheet), cut from a copy of 60.pdf fetched
% through the project's relay
% (https://grothendieck-relay.onrender.com/source/60.pdf); tiles in
% archives/tiles/60/p21-p40, with tighter pdftoppm crops for doubtful words.
% Per the skill transcribe-grothendieck. The folder is sixty-nine pages in
% four batches (1-20, 21-40, 41-60, 61-69), transcribed in parallel; this is
% the second. Batch 1 (transcripts/60/batch-01.fr.tex) was read first and its
% notation is kept.
%
% Pass: Opus 5.5 (claude-opus-5-5), 2026-09-26 — first pass, unchecked against the pages by a human.
%
% Pages skipped: 21, 23, 25, 27, 29, 31, 33, 35, 37, 39 — the typed versos of
% the sheets (Bourbaki material: pages of rédaction n° 392 and the
% « Compléments au chapitre VII d'Algèbre », rédaction n° 397, one of them
% scanned upside down), with nothing in his hand: skipped without
% description by the settled rule.
%
% Pages summarised: none. Pages transcribed: 22, 24, 26, 28, 30, 32, 34, 36,
% 38, 40 — ten rectos in ink, on the backs of the typescripts. His own
% pagination: none visible (the pencilled numbers bottom left are the
% archivists'). No numbered formulas in this batch: the run (1)-(3) of
% batch 1 is not continued; his boxed formulas are kept boxed. Fast drafting
% on pp. 36-40: formulas carry, much of the connective prose does not.
%
% What the batch is about. It continues batch 1 (functional equation of
% L-functions over F_p, delta(E), chi(E), the correction factor A(t) at
% infinity). P. 22: for a « motif » E, chi(E) and delta(E) in terms of the
% virtual Betti numbers beta_i(E) and the virtual components h^i_!(E),
% delta(h^i_!) = eps_i p^{i beta_i/2}, whence delta(E) = eps(E)
% p^{(1/2) sum (-1)^i i beta_i}. P. 24: the term at infinity
% R_infty f(E) = R f_* j^*(R i_*(E)); by local duality
% D R_infty f(E) = R_infty f(D E)[1] (« 1 ou -1 ! »), and for E lisse of
% weight rho on X smooth of dim n, the functional equation of L^infty_E:
% L^infty_E(1/(p^{n+rho} t)) L^infty_E(t) (-t)^{chi_infty} delta_infty = 1.
% P. 26: chi_infty(E) = 0, delta_infty(E) = p^{(n+rho)chi} delta^{-2} =
% eps(E)^{-2}; a boxed list of « invariants arithmétiques intéressants »
% delta, delta', eps, delta_infty with q = p^{n+rho}. Pp. 28-30: consistency
% check: B(t) = A(t) delta(E) (-t)^chi satisfies B(t)B(1/(qt)) = 1, which is
% again the equation of L^infty; with xi_E(t) = t^{chi/2} L_E(t),
% xi_E(1/(qt)) = A(t) eps'(E) xi_E(t), eps'(E) = eps(E)(-1)^chi; X complete:
% A = 1, eps' = ±1. Pp. 32-34: the case of a curve: X smooth projective
% over F_p with function field K, E_eta a Q_l-sheaf on eta = Spec K,
% L^*_{E_eta} := L_{i_* E_eta}, D(i_* E_eta) = i_*(E_eta^v)(1), the
% functional equation of L^* and of xi^* with q = p^{1+rho}; remark that L^*
% is not additive in E_eta: a correction by prod_x P_{Q_x}(t) at the points
% of ramification. Pp. 36-40: attempts to make it additive (via composition
% series, not local in (E,x)); a group D, a normal subgroup I with an open
% subgroup U acting unipotently; the « moyenne » E -> E^I from Mod(D/U,k) to
% Mod(D/I,k), exact in characteristic 0, inducing homomorphisms of
% representation rings R(D/U,k) -> R(D/I,k) compatible for V ⊂ U (a small
% triangle, drawn with tikzcd).
%
% Notation, as in batch 1: f for Frobenius; \check E for the contragredient;
% D for the dualising functor on p. 22-34, but on pp. 38-40 D is a group
% (décomposition) — both are his; R f_*, R f_!, R_\infty f with plain R; he
% writes R_\infty f and R^*_\infty f on p. 24 and both are kept; L^\infty_E;
% \varepsilon for the Euler characteristic of a complex on p. 26, as on
% batch 1 p. 20; \mathcal-free E_\eta, \eta = Spec K. The map symbol shaped
% like a lightning flash on p. 40 is rendered \natural_U, with a note.
%
% Readings to check first: « \uncertain{toujours} entier » on p. 22; the top
% line and the right-hand prose of p. 24; the last box of p. 26 (the sign of
% the exponent of q in eps(E) = delta'(E) q^{-chi/2}); the superscript on
% A(t) on p. 30; the looped insertion on p. 34; the heading and margin NB
% of p. 38; nearly all the prose of pp. 36 and 40.

\documentclass[11pt,a4paper]{article}
% Le préambule commun à toutes les transcriptions du fonds Grothendieck.
%
% Il tient en une idée : **l'incertitude se compose, elle ne se résout pas.**
% Une transcription qui lisse un mot douteux a détruit la seule chose pour
% laquelle on la faisait — un lecteur doit pouvoir savoir, à chaque endroit,
% ce qui était sur le feuillet et ce qui a été ajouté. Les six macros qui
% suivent sont donc l'appareil critique, et non de la décoration.
%
% Les mêmes commandes sont reconnues par `scripts/render.mjs`, qui produit la
% vue de lecture du site. Une macro ajoutée ici doit l'être aussi là-bas,
% faute de quoi le rendu échouera bruyamment — ce qui est voulu : un rendu
% silencieusement faux est pire qu'un rendu absent.

\ProvidesPackage{grothendieck}[2026/08/08 Transcriptions du fonds Grothendieck]

\usepackage{amsmath,amssymb,amsthm}
\usepackage{mathrsfs}
\usepackage{stmaryrd}
% Les diagrammes commutatifs sont partout dans ce fonds : c'est la langue
% même de la géométrie algébrique de Grothendieck, et une transcription qui
% les rendrait en prose ne transcrirait rien.
\usepackage{tikz-cd}
% Français par défaut — c'est la langue des feuillets — mais l'anglais est
% chargé aussi : la traduction et le résumé sont des documents anglais, et la
% typographie française (espaces avant les deux-points) y serait une faute.
% Ces documents basculent par \selectlanguage{english} en tête de corps.
\usepackage[english,french]{babel}
\usepackage{fontspec}
\usepackage{geometry}
\geometry{a4paper,margin=25mm,marginparwidth=28mm}
\usepackage{marginnote}
\usepackage{xcolor}
% ulem, not soul. `soul`'s \st and \ul are built on a character-by-character
% scan that XeLaTeX's font machinery breaks: the document fails with an
% undefined control sequence rather than degrading. ulem does the same two jobs
% and is Unicode-engine safe. `normalem` stops it hijacking \emph, which the
% transcriptions use for Grothendieck's own emphasis.
\usepackage[normalem]{ulem}
\usepackage{hyperref}
\hypersetup{colorlinks=true,linkcolor=black,urlcolor=black,pdfborder={0 0 0}}

% --- Métadonnées du lot -----------------------------------------------------
% Reprises telles quelles par le script de rendu, qui les affiche en tête de
% la vue de lecture. Elles fixent ce que le fichier prétend couvrir : sans
% elles, une transcription flotte sans point d'attache dans un fonds de seize
% mille feuillets.

\newcommand{\folder}[1]{\gdef\grothFolder{#1}}
\newcommand{\batch}[1]{\gdef\grothBatch{#1}}
\newcommand{\leaves}[2]{\gdef\grothFirst{#1}\gdef\grothLast{#2}}
\newcommand{\foldertitle}[1]{\gdef\grothFoldertitle{#1}}
\gdef\grothFolder{?}\gdef\grothBatch{?}\gdef\grothFirst{?}\gdef\grothLast{?}\gdef\grothFoldertitle{}

% --- Le résumé --------------------------------------------------------------
%
% Il ouvre la lecture modernisée, sous le titre « Résumé », et il est composé
% en petit. Ce n'est pas un ornement : c'est le seul passage du document qui
% ne soit pas des mathématiques, et un lecteur doit pouvoir le reconnaître
% avant de l'avoir lu — puis le sauter s'il connaît déjà le sujet, ou s'y
% arrêter s'il ne le connaît pas. Le filet qui le suit marque la reprise du
% corps.

\newenvironment{resume}
  {\small\setlength{\parskip}{0.45em}}
  {\par\medskip\hrule\medskip}

% --- Mots-clés ---------------------------------------------------------------
%
% En anglais, en clôture du résumé de la lecture modernisée : le vocabulaire
% moderne sous lequel chercher ce que les feuillets construisent. Le manifeste
% du site (`npm run manifest`) les extrait pour étiqueter la cote entière —
% cette ligne est la source unique des tags, il n'y a pas de fichier à côté.
\newcommand{\keywords}[1]{\par\smallskip\noindent
  {\footnotesize\textsc{Keywords} --- \textit{#1}}\par}

% --- L'appareil critique ----------------------------------------------------

% Le numéro de feuillet, en marge. C'est l'ancre qui permet de régler un
% désaccord sur une formule en regardant *un* feuillet plutôt qu'une chemise
% de six cents. Le site s'en sert aussi pour tourner les pages du fac-similé
% quand on fait défiler la transcription.
\newcommand{\leaf}[1]{%
  \marginnote{\footnotesize\color{gray!70}\textsf{#1}}%
  \label{leaf:#1}\ignorespaces}

% Illisible : on ne devine pas. Un mot inventé qui se lit comme les autres est
% le pire résultat possible.
\newcommand{\ill}{\textcolor{red!60!black}{[\,\ldots\,]}}

% Lecture douteuse : le mot est donné, et signalé comme incertain.
\newcommand{\uncertain}[1]{\uline{#1}}

% Ajout de l'éditeur — une lettre manquante, un mot sous-entendu.
\newcommand{\add}[1]{\textcolor{blue!50!black}{[#1]}}

% Biffé par Grothendieck lui-même. Ce qu'il a rayé fait partie du document :
% une rature dit souvent où la pensée a changé de direction.
\newcommand{\struck}[1]{\sout{#1}}

% Note du transcripteur, sur l'état matériel du feuillet.
\newcommand{\note}[1]{\par\smallskip{\small\color{gray!60!black}\textit{[#1]}}\par\smallskip}

% Note marginale de Grothendieck, distincte du corps du texte.
\newcommand{\marginal}[1]{\par\smallskip\noindent\hspace{1em}%
  {\small\color{gray!60!black}#1}\par\smallskip}

% --- Titre ------------------------------------------------------------------

\AtBeginDocument{%
  \begin{center}
    {\large\bfseries Fonds Alexandre Grothendieck --- cote n\textsuperscript{o}~\grothFolder}\\[2pt]
    {\itshape\grothFoldertitle}\\[4pt]
    {\small feuillets \grothFirst--\grothLast\ (lot \grothBatch)}\\[2pt]
    {\footnotesize\color{gray!60!black}Transcription établie d'après le fac-similé numérisé
     par l'Université de Montpellier.\\
     Les lectures incertaines sont soulignées, les illisibles notées [\,\ldots\,],
     les ajouts entre crochets.}
  \end{center}
  \vspace{1em}}

\endinput


\folder{60}
\batch{2}
\pages{21}{40}
\dating{[à partir de 1962- à partir de 1971]}
\watermark{Édition de démonstration}
\foldertitle{Fonctions L / Équation fonctionnelle des fonctions L (cas géométrique) : notes manuscrites (s.d.)}

\begin{document}

\page{22}
\emph{Dans le cas d'un motif $E$},
\note{ces mots sont entourés d'un trait en haut de la page ; « motif » est souligné}
on va exprimer $\delta(E)$ et $\chi(E)$ en termes des \struck{\ill{}}
composantes virtuelles des $R f_{!}(E)$ (utilisant les conj.\ de Weil), on
trouve
\[
  \left\lbrace
  \begin{array}{l}
    \chi(E) = \sum (-1)^i \beta_i(E) \\[6pt]
    \delta(E) = \prod_i \delta\bigl(h^i_{!}(E)\bigr)^{(-1)^i}
  \end{array}
  \right.
\]
\marginal{en face : $\beta_i(E)$ les nbs de Betti virtuels à coeff.\ dans $E$}
Or $h^i_{!}(E)$ donne \uncertain{des} valeurs propres de val.\ absolue
$p^{i/2}$, donc
\[
  \delta\bigl(h^i_{!}(E)\bigr) = \varepsilon_i\, p^{\frac{i \beta_i}{2}}
\]
(NB. $i\beta_i/2$ \uncertain{toujours} entier)
\note{devant $\delta(h^i_{!}(E))$, un premier départ est biffé}
où \struck{\ill{}} $\varepsilon_i$ est défini comme un signe de déterminant
virtuel, et $\varepsilon_i$ est $= +1$ si $i$ est impair. On trouve donc
\[
  \left\lbrace
  \begin{array}{l}
    \delta(E) = \varepsilon(E)\, p^{\frac{1}{2}\sum_i (-1)^i i\, \beta_i(E)} \\[6pt]
    \varepsilon(E) = \prod_i \varepsilon_i(E)
  \end{array}
  \right.
\]
\note{entre $\delta(E) =$ et $\varepsilon(E)$, un premier membre de droite,
une puissance de $p$, est lourdement biffé}
\marginal{en face : $\varepsilon(E)$ et les $\varepsilon_i(E)$ sont $\pm 1$,
$\varepsilon_i(E) = +1$ si $i$ impair.}

\page{24}
\ill{} \ill{} de $L_E$ et $L^{\infty}_{DE}$ \ldots
\struck{$R f_{*}\, j^{*} R i_{*}(E)$}
\[
  R_{\infty} f(E) = R f_{*}\, j^{*}\bigl(R i_{*}(E)\bigr)
\]
\[
  \begin{aligned}
    D\, R_{\infty} f(E) &= D\, R f_{*}(\quad) \\
      &= R f_{*}\, D(\quad) \\
      &= R f_{*}\, R j^{!}\bigl(D\, R i_{*}(E)\bigr) \\
      &= R f_{*}\Bigl(R j^{!}\bigl(R i_{!}(D E)\bigr)\Bigr)
  \end{aligned}
\]
\marginal{une accolade embrasse les deuxième et troisième lignes : \emph{dualité locale}}
\note{les parenthèses des deux premières lignes sont laissées vides ; sous
la dernière, une accolade renvoie à $j^{*}\bigl(R i_{*}(D(E))\bigr)[1]$}
\[
  \boxed{D\, R^{*}_{\infty} f(E) = R^{*}_{\infty} f\bigl(D(E)\bigr)\, [1]}
\]
\marginal{en face : \ill{} \ill{} \ill{} $1$ ou $-1$ !}
\note{une ligne suivante, $L^{\infty}_E \ldots L^{\infty}_{\check E} \ldots = 1$,
est biffée de plusieurs traits}

$L^{\infty}_{D(E)}$ \struck{Supposons \ldots\ $D(E) \simeq \check E(n+\rho)$,
on trouve}

Supposons $E$ loc.\ ct.\ \add{de poids $\rho$} et $X$ lisse de dim.\ \ill{}
\ill{}, \ill{} \ill{} \ill{} :
\note{« de poids $\rho$ » est écrit au-dessus de la ligne ; la fin de la
phrase, à droite, se lit mal}
\[
  D(E) \simeq \check E(n) \simeq E(n + \rho)
\]
d'où
\[
  \boxed{D\, R^{*}_{\infty} f(E) \simeq R^{*}_{\infty} f(E)\,(n + \rho)\,[1]}
\]
donc \ill{} l'équation fonctionnelle
\[
  (-t)^{-\chi_{\infty}(E)}\, \delta_{\infty}(E)\, L^{\infty}_E(t^{-1})
  = L^{\infty}_E\bigl(p^{-(n+\rho)} t\bigr)^{-1}
\]
\note{devant $(-t)$, un signe biffé}
i.e.
\[
  \boxed{L^{\infty}_E\Bigl(\frac{1}{p^{n+\rho}\, t}\Bigr) L^{\infty}_E(t)\,
    (-t)^{\chi_{\infty}(E)}\, \delta_{\infty}(E) = 1}
\]

\page{26}
\emph{NB}
\[
  \chi_{\infty}(E) = \varepsilon\bigl(R f_{*}(E)\bigr) - \varepsilon\bigl(R f_{!}(E)\bigr)
\]
or on a vu que les deux $\varepsilon$ sont égaux, donc
$\chi_{\infty}(E) = 0$. \uncertain{D'autre part}
\note{comme page 20, $\varepsilon$ note ici la caractéristique d'Euler du
complexe}
\[
  \delta_{\infty}(E) = \delta\bigl(R f_{*}(E)\bigr)\, \delta\bigl(R f_{!}(E)\bigr)^{-1}
\]
et on a vu que
\[
  p^{(n+\rho)\chi(E)}\, \delta\bigl(R f_{*}(E)\bigr)^{-1} = \delta\bigl(R f_{!}(E)\bigr)
\]
d'où
\[
  \left\lbrace
  \begin{array}{l}
    \delta_{\infty}(E) = p^{(n+\rho)\chi(E)}\, \delta(E)^{-2} = \varepsilon(E)^{-2} \\[6pt]
    \varepsilon(E) = \delta(E)\, p^{-(n+\rho)\chi(E)/2}
  \end{array}
  \right.
\]
\note{après $\varepsilon(E)^{-2}$, une expression est lourdement biffée ; le
$\varepsilon$ de ces deux lignes est tracé en gras}
($= \pm 1$ si $X$ est \struck{\ill{}} propre)

(qui est bien $1$ lorsque $X$ est complet !)

d'où l'équation fonctionnelle
\[
  \boxed{L^{\infty}_E\Bigl(\frac{1}{p^{n+\rho}\, t}\Bigr) L^{\infty}_E(t)\,
    \delta_{\infty}(E) = 1}
\]
\[
  \Bigl(\delta_{\infty}(E) = p^{(n+\rho)\chi(E)}\, \delta(E)^{-2} = \varepsilon(E)^{-2}\Bigr)
\]
\note{avant le dernier signe $=$, un symbole est biffé}

[Invariants arithmétiques \uncertain{intéressants}, (moyennant la
connaissance de l'invariant géométrique $\chi(E)$, et de $q = p^{n+\rho}$) :
\[
  \delta(E) = \delta^{*}_{!}(E) , \quad
  \delta'(E) = \delta^{*}_{*}(E) , \quad
  \varepsilon(E) , \quad
  \delta_{\infty}(E) ]
\]
\note{dans $q = p^{n+\rho}$, le $q$ surcharge une autre lettre}
\[
  \boxed{\delta(E)\,\delta'(E) = q^{\chi(E)}}
  \qquad
  \boxed{\varepsilon(E) = \delta(E)\, q^{-\chi(E)/2} = \delta'(E)\, q^{-\chi(E)/2}}
\]
\[
  \boxed{\delta_{\infty}(E) = q^{\chi(E)}\, \delta(E)^{-2} = q^{-\chi(E)}\, \delta'(E)^{2}
    = \varepsilon(E)^{-2}}
\]
\marginal{une flèche part du signe $=$ de la deuxième formule encadrée :
variance avec $E$, $X$ !!}
\marginal{dans un ovale, à gauche : NB.\ $\varepsilon(E) = \pm 1$ si $X$ complet}
\note{la page écrit bien $\delta'(E)\, q^{-\chi(E)/2}$ dans la deuxième
formule encadrée ; d'après la première, on attendrait
$\delta'(E)^{-1} q^{\chi(E)/2}$. On transcrit tel quel. L'exposant $2$ de
$\delta'(E)$ dans la troisième formule, au bord de la feuille, est douteux}

\page{28}
Comparons cette équation fonctionnelle avec celle pour la fonction $L_E$, on
trouve, en \struck{substituant} posant pour simplifier $p^{n+\rho} = q$, et
$A(t)\, \delta(E)\, (-t)^{\chi(E)} = B(t)$ :
\[
  L\Bigl(\frac{1}{qt}\Bigr) = B(t)\, L(t)
\]
et faisant $t \mapsto \frac{1}{qt}$ dedans, on trouve
\[
  L(t) = B\Bigl(\frac{1}{qt}\Bigr) L\Bigl(\frac{1}{qt}\Bigr)
  = B\Bigl(\frac{1}{qt}\Bigr) B(t)\, L(t)
\]
d'où
\[
  B(t)\, B\Bigl(\frac{1}{qt}\Bigr) = 1 .
\]
Remplaçant $B$ par sa valeur, on trouve
\[
  L^{\infty}_E\Bigl(\frac{1}{qt}\Bigr)^{-1} \delta(E)\, (-t)^{\chi(E)}\,
  L^{\infty}_E(t)^{-1}\, \delta(E) \Bigl(\frac{1}{qt}\Bigr)^{\chi(E)} = 1
\]
\note{dans le premier argument, un $p$ biffé précède $qt$ ; le dernier
facteur est surchargé et pourrait porter un signe $-$}
i.e.
\[
  L^{\infty}_E(t)\, L^{\infty}_E\Bigl(\frac{1}{qt}\Bigr) = \delta(E)^2\, q^{-\chi(E)}
\]
\note{devant $q^{-\chi(E)}$, un petit symbole biffé ; au-dessous, une ligne
commençant par $L^{\infty}_E\bigl(\frac{1}{p^{n+\rho} t}\bigr)$ est biffée}
C'est précisément l'équation fonctionnelle explicitée plus haut.

Posons
\[
  \xi_E(t) = t^{\frac{\chi(E)}{2}} L_E(t)
\]
alors
\[
  \begin{aligned}
    \xi_E\Bigl(\frac{1}{qt}\Bigr)
      &= \Bigl(\frac{1}{qt}\Bigr)^{\chi(E)/2} L_E\Bigl(\frac{1}{qt}\Bigr)
       = \Bigl(\frac{1}{qt}\Bigr)^{\chi(E)/2} A(t)\, \delta(E)\, (-t)^{\chi(E)} L_E(t) \\
      &= q^{-\chi(E)/2} A(t)\, \delta(E)\, (-1)^{\chi(E)}\; t^{\chi(E)/2} L_E(t)
  \end{aligned}
\]
i.e.
\note{sous la dernière ligne, deux accolades : l'une sous
$q^{-\chi(E)/2} \ldots (-1)^{\chi(E)}$, l'autre sous $t^{\chi(E)/2} L_E(t)$,
légendée $\xi_E(t)$}

\page{30}
\[
  \boxed{\xi_E\Bigl(\frac{1}{qt}\Bigr) = A(t)\, \varepsilon'(E)\, \xi_E(t)}
  \qquad q = p^{n\rho}
\]
\note{$A$ porte un petit exposant qui ne se lit pas ; la page écrit
$q = p^{n\rho}$, sans signe $+$ visible, là où les pages 26-28 ont
$q = p^{n+\rho}$}
\[
  \left\lbrace
  \begin{array}{l}
    A(t) = L^{\infty}_E\Bigl(\dfrac{1}{qt}\Bigr)^{-1} \\[8pt]
    \varepsilon'(E) = \varepsilon(E)\, (-1)^{\chi(E)}
  \end{array}
  \right.
  \qquad
  \varepsilon(E) = \delta(E)\, q^{-\chi(E)/2}
\]
\note{dans la dernière formule, un premier membre de droite est biffé}

Si \struck{\ill{}} $X$ complète
$\left\lbrace
\begin{array}{l}
  A(t) = 1 \\
  \varepsilon'(E) = \pm 1
\end{array}
\right.$

\page{32}
Cas d'une courbe alg.

$X$ \add{\uncertain{courbe},} lisse projective [\uncertain{irréd.}] sur
$\mathbb{F}_p$, \struck{\ill{}} corps de fonctions $K$.
\note{le mot entre virgules et le crochet sont écrits au-dessus de la ligne}

$E_\eta$ $\mathbb{Q}_\ell$-faisceau constr.\ sur $\eta = \operatorname{Spec} K$.

On veut lui associer une « fonction $L$ ». C'est \uncertain{raisonnable},
car $K$ détermine $X$.

Prenons d'abord $E = i_{*}(E_\eta)$, où $i : \eta \to X$.

On sait alors que
\[
  D(E) = i_{*}(\check E_\eta)(1)
\]
\struck{Donc si on suppose} Donc si on pose
$L^{*}_{E_\eta}(t) = L_{i_{*}(E_\eta)}(t)$, on aura
\[
  L^{*}_{\check E_\eta}\Bigl(\frac{1}{pt}\Bigr)
  = (-t)^{\chi^{*}(\check E_\eta)}\, \delta^{*}(E_\eta)\, L^{*}_{E_\eta}(t)
\]
\note{devant $\bigl(\frac{1}{pt}\bigr)$, un symbole biffé}
si on pose
\[
  \chi^{*}(E_\eta) = \chi\bigl(i_{*}(E_\eta)\bigr) , \qquad
  \delta^{*}(E_\eta) = \delta\bigl(i_{*}(E_\eta)\bigr)
\]
Supposons qu'on ait
\[
  \check E_\eta = E_\eta(\rho) \qquad \text{d'où} \qquad
  i_{*}(\check E_\eta) = i_{*}(E_\eta)(\rho)
\]
alors on trouve,
\[
  \boxed{L^{*}_{E_\eta}\Bigl(\frac{1}{qt}\Bigr)
    = (-t)^{\chi^{*}(E_\eta)}\, \delta^{*}(E_\eta)\, L^{*}_{E_\eta}(t)}
  \qquad (q = p^{1+\rho}
\]

\page{34}
Posant
\[
  \xi^{*}_{E_\eta}(t) = t^{\chi^{*}(E_\eta)/2} L^{*}_{E_\eta}(t)
\]
on trouve
\[
  \boxed{\xi^{*}_{E_\eta}\Bigl(\frac{1}{qt}\Bigr)
    = \varepsilon'^{*}(E_\eta)\, \xi^{*}_{E_\eta}(t)}
\]
\[
  \left\lbrace
  \begin{array}{l}
    \varepsilon'^{*}(E_\eta) = \delta^{*}(E_\eta)\, q^{-\chi^{*}(E_\eta)/2}\, (-1)^{\chi^{*}(E_\eta)} \\[6pt]
    q = p^{1+\rho}
  \end{array}
  \right.
\]
\marginal{en face : \ill{} \ill{} \ill{}}

\emph{Une remarque} : \uncertain{fonction} $L^{*}_{E_\eta}(t)$ :
\emph{\ill{} \ill{} \ill{}} \emph{additif} en $E$. De façon précise, si on a
\[
  0 \to E'_\eta \to E_\eta \to E''_\eta \to 0
\]
on en conclut
\[
  0 \to i_{*}(E'_\eta) \to i_{*}(E_\eta) \to i_{*}(E''_\eta) \to Q \to 0
\]
où
\struck{$Q = \operatorname{Ker}\bigl(R^1 i_{*}(E'_\eta) \to R^1 i_{*}(E_\eta)\bigr)$}
\[
  Q \simeq \coprod_{x \text{ pt ramifié de } E''_\eta} j_x(\ill{})
\]
\note{la formule biffée est remplacée par l'expression $\coprod$, inscrite
dans une boucle qui renvoie à la fin de la ligne}
est \ill{} \ill{} \ill{} \ill{} de « ramification » de $E''_\eta$.

D'où
\[
  \left\lbrace
  \begin{array}{l}
    L^{*}_{E_\eta}(t) = L^{*}_{E'_\eta}(t)\, L^{*}_{E''_\eta}(t)\, L_Q(t)^{-1}
      = L^{*}_{E'_\eta}(t)\, L^{*}_{E''_\eta}(t) \prod_x P_{Q_x}(t) \\[8pt]
    Q_x = \operatorname{Coker}(E_x \to E''_x)
  \end{array}
  \right.
\]
\marginal{en bas à droite, sous un arc :
$H^2(\bar X, E) =$ dual de $H^0(\bar X, E(\ill{}))(\ill{})$}

\page{36}
On peut songer à la rendre additive, en la définissant d'abord sur les
$E_\eta$ simples (ou \uncertain{abst.}\ simples, quitte à prendre des
extensions finies de $\mathbb{Q}_\ell$) et en prolongeant par linéarité ou
plus \uncertain{simplement} à définir les facteurs locaux dans
$L_{E_\eta}(t)$ par $L_{E_\eta(x)}(t)$, en posant
$E_\eta(x) = \sum E^{i*}_\eta(x)$, les $E^{i}_\eta$ étant les facteurs
simples dans une suite de composition de $E$. Ceci n'est pas encore
satisfaisant, car les $E_\eta(x)$ ainsi construits ne sont pas de nature
locale en $(E, x)$, i.e.\ ne commutent pas à l'extension étale
\struck{sur $X$} (i.e.\ ext.\ résiduellement \add{\ill{}} triviale) sur $X$.
\note{l'indice de $E^{i*}_\eta(x)$ est douteux ; au-dessus de
« triviale », un mot ajouté, puis un mot biffé}

\page{38}
\section{Suite de la \uncertain{remarque}}
\note{titre souligné, de sa main, en haut de la page}
\marginal{dans un cadre, en haut à droite : NB. S'il s'agit de \ill{} des
\ill{} galoisiens, et $U$ est \uncertain{ouvert} dans $I$, alors il
s'ensuit qu'il contient un s-gpe ouvert \uncertain{de} $U$, qui est
\emph{invariant} par $D$}

$D$ groupe, $I$ sous-groupe distingué, $E$ vectoriel de dim.\ finie sur un
corps $k$, sur lequel $D$ opère. On suppose que $I$ admet un sous-groupe
d'indice fini $U$, \add{invariant par $D$,} tel que pour $g \in U$, $g$
opère sur $E$ de façon \emph{unipotente}, i.e.\ valeurs propres toutes $1$.
Alors le groupe algébrique d'automorphismes de $E$ engendré par les
$g \in U$ est « unipotent », donc \struck{\ill{}} $E$ admet une suite de
composition sous $U$ telle que \add{\uncertain{$E$ quot.}} \ill{} \ill{}
facteurs. \struck{soit stable par $E$.} De façon précise, on peut prendre la
suite de composition formée par les $E^U$, $(E/E^U)^U$, \ldots. Cette suite
est stable par les opérations de $D$, si \struck{(comme il est loisible)}
\struck{on suppose} \add{puisque} $U$ dans $I$ invariant par les $D$. Donc il
existe une suite de composition de $E$ \emph{sous $D$} telle que sur chaque
facteur $E_i$, \struck{\ill{}} $U$
\note{« invariant par $D$ » est ajouté au-dessus de la ligne ; la fin de la
proposition sur les facteurs, entourée d'une boucle, se lit mal ; la
phrase continue page 40}

\page{40}
opère trivialement, \add{(et \uncertain{respectant} bien \ill{})} \ill{}
\ill{} \ill{} \ill{} au virtuel, \add{et pour $U$ fixé,} les modules $E$
envisagés sont assimilables à ceux pour lesquels \struck{il existe un s-gpe
\ill{}} $U$ opère trivialement sur $E$.
\note{les deux ajouts sont écrits au-dessus de la ligne}

\ill{} \ill{} $E$ sur lequel $U$ opère trivialement, on peut prendre
l'opération moyenne
\[
  E \mapsto E^{I} \quad \text{de} \quad
  \mathrm{Mod}(D/U, k) \to \mathrm{Mod}(D/I, k)
\]
et si $k$ est de car.\ $0$, c'est le \uncertain{seul} foncteur exact. Il
induit un \ill{} \struck{\ill{} opération additive, \ill{} \ill{} \ill{}}
\struck{\ill{}} naturel \struck{\ill{}} $E \mapsto E^{I}$
\[
  \natural_U : R(D/U, k) \to R(D/I, k)
\]
\note{le symbole de l'application a la forme d'un éclair ; on le rend
conventionnellement par $\natural$}
\uncertain{On a} compatibilité, pour $V \subset U$, par les foncteurs
$E \mapsto E^{I}$, d'où par les $\natural_U$ et $\natural_V$ :

\begin{tikzcd}
R(D/U, k) \arrow[r, "\natural_U"] \arrow[d] & R(D/I, k) \\
R(D/V, k) \arrow[ur, "\natural_V"'] &
\end{tikzcd}

\note{la flèche verticale porte une légende de deux mots qui ne se lisent
pas}

\end{document}
